% from ega0-20-fr.tex
\oldpage[0\textsubscript{IV-1}]{117}
\begin{proposition}[20.1.1]
\label{0.20.1.1-fr}
Soient $A$ un anneau (non nécessairement commutatif), $B$, $E$, $C$ des
$A$-anneaux, $p:E\to C$ un $A$-homomorphisme dont le noyau $\mathfrak{Q}$
est de carré nul, $u:B\to C$ un $A$-homomorphisme. Supposons qu'il existe
un $A$-homomorphisme $v_0:B\to E$ tel que $u$ se factorise en
$B\xrightarrow{v_0}E\xrightarrow{p}C$. Alors l'ensemble des
$A$-homomorphismes $v:B\to E$ tels que $u=p\circ v$ est identique à
l'ensemble des applications $v_0+D$, où $D:B\to\mathfrak{Q}$ est une
application vérifiant les deux conditions suivantes :
\begin{enumerate}
  \item[\rm(i)] $D$ est un homomorphisme de $A$-bimodules.
  \item[\rm(ii)] Quels que soient $f$, $g$ dans $B$, on a
\end{enumerate}
\begin{equation}
\tag{20.1.1.1}
\label{0.20.1.1.1-fr}
  D(fg)=f.D(g)+D(f).g.
\end{equation}
\end{proposition}

\begin{proof}
Dire que $p(v(f))=p(v_0(f))$ pour $f\in B$ signifie que
$D(f)=v(f)-v_0(f)$ appartient à $\mathfrak{Q}$ ; si l'on écrit que
$v(fg)=v(f)v(g)$, on obtient la relation \sref{0.20.1.1.1-fr},
$\mathfrak{Q}$ étant de carré nul, et la condition (i) résulte de ce que $v$
et $v_0$ sont des $A$-homomorphismes.

Si $\rho:A\to B$ est l'homomorphisme structural, on tire de
\sref{0.20.1.1.1-fr} que l'on a
$D(\rho(a)f)=\rho(a)D(f)+D(\rho(a))f$ pour tout $a\in A$ ; mais on doit
avoir $D(\rho(a)f)=\rho(a)D(f)$ d'après (i), donc, en faisant $f=1$, il
vient
\[
  D(\rho(a))=0 \qquad\text{pour}\quad a\in A\,;
\]
réciproquement, si $D$ est nulle dans $\rho(A)$ et vérifie
\sref{0.20.1.1.1-fr}, elle vérifie aussi (i).
\end{proof}

\begin{definition}[20.1.2]
\label{0.20.1.2-fr}
Étant donnés un $A$-anneau $B$ et un $B$-bimodule $L$, on appelle
\emph{$A$-dérivation de $B$ dans $L$} une application $D:B\to L$ vérifiant
les conditions (i) et (ii) de \sref{0.20.1.1-fr}.
\end{definition}

Il résulte de \sref{0.20.1.1-fr} que le \emph{noyau} d'une $A$-dérivation
$D:B\to L$ est un sous-$A$-anneau de $B$.

On appelle parfois \emph{dérivation} de $B$ dans $L$ une
$\mathbf{Z}$-dérivation, c'est-à-dire une application \emph{additive} de $B$
dans $L$ vérifiant \sref{0.20.1.1.1-fr} ; on a donc $D(1)=0$ pour une telle
application. Si $B$ est une algèbre sur un \emph{corps premier} $P$, toute
$\mathbf{Z}$-dérivation de $B$ est une $P$-dérivation : c'est évident d'après
ce qui précède si $P$ est de caractéristique $>0$, et dans le cas contraire,
la relation $n.n^{-1}=1$ pour $n\in\mathbf{Z}$ donne
\[
  nD(n^{-1})=0
\]
donc $D(n^{-1})=0$ puisque $P$ est de caractéristique $0$.

Il résulte aussitôt de la définition \sref{0.20.1.2-fr} que si $D$, $D'$ sont
deux $A$-dérivations de $B$ dans $L$, il en est de même de $D-D'$. Autrement
dit, l'ensemble des $A$-dérivations de $B$ dans $L$ est muni d'une structure
de \emph{groupe additif} ; on note ce groupe
$\operatorname{D\acute{e}r}_A(B,L)$. Si $A$ est \emph{commutatif}, $B$ une
$A$-\emph{algèbre commutative}, et $L$ un $B$-module, alors, pour toute
$A$-dérivation $D$ de $B$ dans $L$ et tout $a\in A$, $aD$ est encore une
$A$-dérivation de $B$ dans $L$, autrement dit $\operatorname{D\acute{e}r}_A(B,L)$
est alors muni d'une structure de $A$-\emph{module}.

